\pagestyle{empty}
\tight
\begin{tblr}{width=\textwidth,colspec={|Q[c,m,1.9cm]|X[l,m]|},hlines,vlines,hspan=minimal,row{1}={bg=headblue},rowsep=1pt,colsep=3pt}
\SetCell{c,m}\hfil\logo\hfil & \textbf{POLITEKNIK NEGERI MALANG}\par\textbf{JURUSAN TEKNOLOGI INFORMASI}\par\textbf{PROGRAM STUDI: D4 TEKNIK INFORMATIKA} \\
\SetCell[c=2]{c} \textbf{RENCANA PEMBELAJARAN SEMESTER (RPS)} & \\
\end{tblr}
\begin{longtblr}{width=\textwidth,colspec={|Q[l,m,1.9cm]|X[0.85,l,m]|X[1.45,l,m]|X[0.75,l,m]|X[1.15,l,m]|},hlines,vlines,hspan=minimal,rowsep=1pt,colsep=2pt,row{1,3}={bg=headgray,font=\bfseries}}
MATA KULIAH & KODE & BOBOT (SKS)/jam & SEMESTER & TGL. PENYUSUNAN \\
Proyek Teknologi Terintegrasi & RTI256104 & 6 SKS/12 jam & 6 & 1 Juli 2025 \\
OTORISASI & Dosen Pengembang RPS & Koordinator RMK & \SetCell[c=2]{l} Ka PRODI &  \\
 & Triana Fatmawati, S.T., M.T. & Triana Fatmawati, S.T., M.T. & \SetCell[c=2]{l}{\vspace{8mm}\par Dr. Ely Setyo Astuti, S.T., M.T.} &  \\
\SetCell[r=20]{l} Capaian Pembelajaran (CP) & \SetCell[c=4]{l,bg=headgray,font=\bfseries} Capaian Pembelajaran Lulusan Program Studi (CPL-Prodi) &  &  &  \\
 & CPL02 & \SetCell[c=3]{l} Mampu menerapkan prinsip rekayasa sistem dan perangkat lunak untuk menghasilkan solusi aplikasi multi-platform sesuai kebutuhan. &  &  \\
 & CPL03 & \SetCell[c=3]{l} Mampu mengelola tim, manajemen diri, berkomunikasi lisan/tulisan, dan mempresentasikan dalam konteks profesional teknologi informasi. &  &  \\
 & CPL06 & \SetCell[c=3]{l} Mampu merancang, mengimplementasikan, menguji, melakukan deployment, memelihara, serta menjamin mutu perangkat lunak yang menjawab permasalahan dan memenuhi kebutuhan pengguna. &  &  \\
 & CPL08 & \SetCell[c=3]{l} Mampu mengelola proyek teknologi informasi secara profesional melalui perencanaan, koordinasi, komunikasi, dan optimalisasi sumber daya. &  &  \\
 & \SetCell[c=4]{l,bg=headgray,font=\bfseries} Capaian Pembelajaran mata kuliah (CPL-MK) &  &  &  \\
 & CPMK0801 & \SetCell[c=3]{l} Mahasiswa mampu merencanakan proyek TI terintegrasi dengan estimasi biaya, jadwal, alokasi sumber daya, dan manajemen risiko. &  &  \\
 & CPMK0802 & \SetCell[c=3]{l} Mahasiswa mampu mengimplementasikan siklus hidup rekayasa perangkat lunak secara menyeluruh dalam proyek nyata multi-platform. &  &  \\
 & CPMK0803 & \SetCell[c=3]{l} Mahasiswa mampu mengeksekusi, memonitor, dan mengendalikan proyek serta mengelola dinamika tim secara profesional. &  &  \\
 & CPMK0804 & \SetCell[c=3]{l} Mahasiswa mampu mendemonstrasikan pengelolaan proyek TI utuh dari perencanaan hingga deliverable akhir. &  &  \\
 & CPMK0607 & \SetCell[c=3]{l} Mahasiswa mampu menghasilkan solusi sistem informasi terintegrasi yang menjawab kebutuhan nyata pengguna. &  &  \\
 & CPMK0205 & \SetCell[c=3]{l} Mampu menerapkan prinsip dan pola desain perangkat lunak untuk membuat blueprint arsitektural dan komponen solusi aplikasi multi-platform yang modular, maintainable, dan scalable. &  &  \\
 & CPMK0302 & \SetCell[c=3]{l} Mampu mengelola tim dan proyek secara terencana, mulai dari perencanaan, koordinasi, komunikasi antar anggota, hingga penyampaian hasil kepada pemangku kepentingan. &  &  \\
 & \SetCell[c=4]{l,bg=headgray,font=\bfseries} Kemampuan akhir tiap tahapan belajar (Sub-CPMK) &  &  &  \\
 & SCPMK0801-04801 & \SetCell[c=3]{l} Mahasiswa mampu menyusun dokumen perencanaan proyek TI terintegrasi mencakup WBS, jadwal, anggaran, dan risk register. &  &  \\
 & SCPMK0802-04802 & \SetCell[c=3]{l} Mahasiswa mampu mengimplementasikan aplikasi multi-platform menggunakan metodologi agile dengan sprint yang terdokumentasi. &  &  \\
 & SCPMK0803-04803 & \SetCell[c=3]{l} Mahasiswa mampu melaksanakan monitoring sprint, mengelola backlog, dan menyajikan laporan kemajuan proyek. &  &  \\
 & SCPMK0804-04804 & \SetCell[c=3]{l} Mahasiswa mampu mendemonstrasikan produk akhir proyek TI terintegrasi dengan defense teknis yang komprehensif. &  &  \\
 & SCPMK0607-04805 & \SetCell[c=3]{l} Mahasiswa mampu mengintegrasikan solusi front-end, back-end, dan basis data dalam satu produk yang deployable. &  &  \\
 & SCPMK0205-04806 & \SetCell[c=3]{l} Mahasiswa mampu merancang blueprint arsitektural dan komponen solusi yang modular, maintainable, dan scalable untuk aplikasi multi-platform. &  &  \\
 & SCPMK0302-04807 & \SetCell[c=3]{l} Mahasiswa mampu mengelola dinamika tim proyek dan menyampaikan hasil kerja tim kepada pemangku kepentingan secara profesional. &  &  \\
\SetCell[r=2]{l} Penilaian & \SetCell[c=4]{l} *Nilai CPMK didapat dari agregasi nilai SUB-CPMK yang dibebankan &  &  &  \\
 & \SetCell[c=4]{l}{\tiny\begin{tblr}{width=\linewidth,colspec={|Q[c,m,0.1\linewidth]|Q[c,m,0.16\linewidth]|X[l,m]|X[l,m]|X[l,m]|X[l,m]|X[l,m]|X[l,m]|Q[c,m,0.08\linewidth]|},hlines,vlines,hspan=minimal,rowsep=1pt,colsep=0.7pt,row{1,2}={bg=headgray,font=\bfseries}}
Kode CPMK & Kode SCPMK & \SetCell[c=6]{c} Bobot per Bentuk Penilaian & & & & & & Total Bobot \\
 & & Dok.Perencanaan & Sprint 1-2 & UTS Milestone & Sprint 3-4 & Laporan & UAS Defense & \\
CPMK0801 & SCPMK0801-04801 & 15\%: perencanaan proyek TI & 0\% & 0\% & 0\% & 0\% & 0\% & 15\% \\
CPMK0802 & SCPMK0802-04802 & 0\% & 7\%: sprint implementasi & 0\% & 10\%: sprint implementasi & 0\% & 0\% & 17\% \\
CPMK0205 & SCPMK0205-04806 & 0\% & 3\%: blueprint arsitektural dan tech stack & 0\% & 0\% & 0\% & 0\% & 3\% \\
CPMK0803 & SCPMK0803-04803 & 0\% & 0\% & 15\%: milestone review & 0\% & 0\% & 10\%: defense proyek & 25\% \\
CPMK0804 & SCPMK0804-04804 & 0\% & 0\% & 0\% & 0\% & 9\%: laporan akhir & 8\%: defense proyek & 17\% \\
CPMK0302 & SCPMK0302-04807 & 0\% & 0\% & 0\% & 0\% & 1\%: koordinasi dan komunikasi tim & 2\%: penyampaian hasil ke stakeholder & 3\% \\
CPMK0607 & SCPMK0607-04805 & 0\% & 10\%: integrasi platform & 0\% & 10\%: integrasi platform & 0\% & 0\% & 20\% \\
\SetCell[c=8]{r}\textbf{Total Per Penilaian} & & & & & & & & \textbf{100\%} \\
\end{tblr}} &  &  &  \\
Diskripsi Singkat Mata Kuliah & \SetCell[c=4]{l} Proyek Teknologi Terintegrasi adalah mata kuliah capstone jalur reguler yang mengintegrasikan seluruh kompetensi rekayasa perangkat lunak dan manajemen proyek dalam pengembangan produk nyata multi-platform menggunakan metodologi agile selama satu semester. &  &  &  \\
Bahan Kajian / Materi Pembelajaran & \SetCell[c=4]{l}\begin{enumerate}\item Manajemen proyek TI: perencanaan, WBS, jadwal, anggaran, dan risiko\item Rekayasa perangkat lunak agile: Scrum, sprint planning, backlog, dan velocity\item Pengembangan aplikasi multi-platform: integrasi front-end, back-end, dan basis data\item Monitoring proyek, kontrol kualitas, deployment, dan defense produk akhir\end{enumerate} &  &  &  \\
\SetCell[r=2]{l} Pustaka & \SetCell[c=4]{l} \textbf{Utama:}\begin{enumerate}\item Schwaber, K. 2020. The Scrum Guide. Scrum.org.\item Sommerville, I. 2016. Software Engineering (10th ed.). Pearson.\item PMI. 2021. PMBOK Guide (7th ed.).\end{enumerate} &  &  &  \\
 & \SetCell[c=4]{l} \textbf{Pendukung:} Materi LMS, repositori proyek, dan jobsheet capstone. &  &  &  \\
Media Pembelajaran & \SetCell[c=2]{l} \textbf{Software:} IDE (VS Code/IntelliJ), Git, project management tools (Jira/Trello), cloud platform, LMS. &  & \SetCell[c=2]{l} \textbf{Hardware:} Komputer/laptop, server/cloud, proyektor. &  \\
Nama Dosen Pengampu & \SetCell[c=4]{l} Tim Pengajar Program Studi D4 TI &  &  &  \\
Matakuliah Syarat & \SetCell[c=4]{l} - &  &  &  \\
\end{longtblr}
\begin{longtblr}{width=\textwidth,colspec={|Q[c,m,0.06\textwidth]|X[1.35,l,m]|X[1.08,l,m]|X[1.16,l,m]|X[0.7,l,m]|X[1.24,l,m]|X[1.06,l,m]|X[0.98,l,m]|Q[c,m,0.05\textwidth]|},hlines,vlines,hspan=minimal,rowhead=2,row{1,2}={bg=headgreen,font=\bfseries},rowsep=1pt,colsep=1.5pt}
Mgg.\par Ke & Kemampuan Akhir Yang Direncanakan (Sub-CP-MK) & Bahan kajian (Materi Pembelajaran) & Bentuk dan Metode Pembelajaran & Estimasi Waktu & Pengalaman Belajar Mahasiswa & Kriteria \& Bentuk Penilaian & Indikator Penilaian & Bobot Penilaian (\%) \\
(1) & (2) & (3) & (4) & (5) & (6) & (7) & (8) & (9) \\
1 & SCPMK0801-04801 & Kickoff: pembentukan tim dan requirement gathering. & Modalitas: Blended Learning. Bentuk: Luring/Daring. Metode: Case Method, PBL, praktik proyek, dan presentasi. & 1 x 6 x 50' tatap muka; 1 x 6 x 50' tugas/praktik mandiri. & Mahasiswa mengerjakan aktivitas terarah untuk kickoff dan requirement gathering dan menunjukkan evidence hasil belajar. & Bentuk penilaian: Kickoff dan Requirement Gathering. Mengacu rubrik RTM. & Ketepatan konsep; kualitas implementasi proyek; validasi dan dokumentasi. & 3 \\
2 & SCPMK0801-04801 & Perencanaan: WBS, jadwal, anggaran, risiko. & Modalitas: Blended Learning. Bentuk: Luring/Daring. Metode: Case Method, PBL, praktik proyek, dan presentasi. & 1 x 6 x 50' tatap muka; 1 x 6 x 50' tugas/praktik mandiri. & Mahasiswa mengerjakan aktivitas terarah untuk perencanaan proyek dan menunjukkan evidence hasil belajar. & Bentuk penilaian: Perencanaan Proyek. Mengacu rubrik RTM. & Ketepatan konsep; kualitas implementasi proyek; validasi dan dokumentasi. & 5 \\
3 & SCPMK0802-04802, SCPMK0205-04806 & Arsitektur solusi dan tech stack: blueprint arsitektural yang modular, maintainable, dan scalable. & Modalitas: Blended Learning. Bentuk: Luring/Daring. Metode: Case Method, PBL, praktik proyek, dan presentasi. & 1 x 6 x 50' tatap muka; 1 x 6 x 50' tugas/praktik mandiri. & Mahasiswa mengerjakan aktivitas terarah untuk arsitektur solusi dan tech stack, termasuk merancang blueprint arsitektural yang modular dan scalable, dan menunjukkan evidence hasil belajar. & Bentuk penilaian: Arsitektur Solusi dan Tech Stack. Mengacu rubrik RTM. & Ketepatan konsep; kualitas implementasi proyek; validasi dan dokumentasi; kualitas blueprint arsitektural. & 2 \\
4 & SCPMK0802-04802 & Sprint planning dan setup environment. & Modalitas: Blended Learning. Bentuk: Luring/Daring. Metode: Case Method, PBL, praktik proyek, dan presentasi. & 1 x 6 x 50' tatap muka; 1 x 6 x 50' tugas/praktik mandiri. & Mahasiswa mengerjakan aktivitas terarah untuk sprint planning dan setup environment dan menunjukkan evidence hasil belajar. & Bentuk penilaian: Sprint Planning dan Setup Environment. Mengacu rubrik RTM. & Ketepatan konsep; kualitas implementasi proyek; validasi dan dokumentasi. & 4 \\
5 & SCPMK0802-04802, SCPMK0607-04805 & Sprint 1: core features back-end. & Modalitas: Blended Learning. Bentuk: Luring/Daring. Metode: Case Method, PBL, praktik proyek, dan presentasi. & 1 x 6 x 50' tatap muka; 1 x 6 x 50' tugas/praktik mandiri. & Mahasiswa mengerjakan aktivitas terarah untuk sprint 1 core features back-end dan menunjukkan evidence hasil belajar. & Bentuk penilaian: Sprint 1 Core Features Back-End. Mengacu rubrik RTM. & Ketepatan konsep; kualitas implementasi proyek; validasi dan dokumentasi. & 5 \\
6 & SCPMK0607-04805 & Sprint 1: front-end dan integrasi. & Modalitas: Blended Learning. Bentuk: Luring/Daring. Metode: Case Method, PBL, praktik proyek, dan presentasi. & 1 x 6 x 50' tatap muka; 1 x 6 x 50' tugas/praktik mandiri. & Mahasiswa mengerjakan aktivitas terarah untuk sprint 1 front-end dan integrasi dan menunjukkan evidence hasil belajar. & Bentuk penilaian: Sprint 1 Front-End dan Integrasi. Mengacu rubrik RTM. & Ketepatan konsep; kualitas implementasi proyek; validasi dan dokumentasi. & 5 \\
7 & SCPMK0803-04803 & Sprint 1 review dan retrospective. & Modalitas: Blended Learning. Bentuk: Luring/Daring. Metode: Case Method, PBL, praktik proyek, dan presentasi. & 1 x 6 x 50' tatap muka; 1 x 6 x 50' tugas/praktik mandiri. & Mahasiswa mengerjakan aktivitas terarah untuk sprint 1 review dan retrospective dan menunjukkan evidence hasil belajar. & Bentuk penilaian: Sprint 1 Review dan Retrospective. Mengacu rubrik RTM. & Ketepatan konsep; kualitas implementasi proyek; validasi dan dokumentasi. & 5 \\
8 & SCPMK0803-04803 & UTS Milestone review dan demo Sprint 1-2. & Modalitas: Blended Learning. Bentuk: Luring/Daring. Metode: Case Method, PBL, praktik proyek, dan presentasi. & 1 x 6 x 50' tatap muka; 1 x 6 x 50' tugas/praktik mandiri. & Mahasiswa mengerjakan aktivitas terarah untuk UTS milestone review dan demo dan menunjukkan evidence hasil belajar. & Bentuk penilaian: UTS Milestone Review dan Demo Sprint 1-2. Mengacu rubrik RTM. & Ketepatan konsep; kualitas implementasi proyek; validasi dan dokumentasi. & 15 \\
9 & SCPMK0607-04805 & Sprint 2: integrasi modul dan database. & Modalitas: Blended Learning. Bentuk: Luring/Daring. Metode: Case Method, PBL, praktik proyek, dan presentasi. & 1 x 6 x 50' tatap muka; 1 x 6 x 50' tugas/praktik mandiri. & Mahasiswa mengerjakan aktivitas terarah untuk sprint 2 integrasi modul dan database dan menunjukkan evidence hasil belajar. & Bentuk penilaian: Sprint 2 Integrasi Modul dan Database. Mengacu rubrik RTM. & Ketepatan konsep; kualitas implementasi proyek; validasi dan dokumentasi. & 4 \\
10 & SCPMK0802-04802 & Sprint 2: testing dan debugging. & Modalitas: Blended Learning. Bentuk: Luring/Daring. Metode: Case Method, PBL, praktik proyek, dan presentasi. & 1 x 6 x 50' tatap muka; 1 x 6 x 50' tugas/praktik mandiri. & Mahasiswa mengerjakan aktivitas terarah untuk sprint 2 testing dan debugging dan menunjukkan evidence hasil belajar. & Bentuk penilaian: Sprint 2 Testing dan Debugging. Mengacu rubrik RTM. & Ketepatan konsep; kualitas implementasi proyek; validasi dan dokumentasi. & 5 \\
11 & SCPMK0803-04803 & Sprint 2 review dan Sprint 3 planning. & Modalitas: Blended Learning. Bentuk: Luring/Daring. Metode: Case Method, PBL, praktik proyek, dan presentasi. & 1 x 6 x 50' tatap muka; 1 x 6 x 50' tugas/praktik mandiri. & Mahasiswa mengerjakan aktivitas terarah untuk sprint 2 review dan sprint 3 planning dan menunjukkan evidence hasil belajar. & Bentuk penilaian: Sprint 2 Review dan Sprint 3 Planning. Mengacu rubrik RTM. & Ketepatan konsep; kualitas implementasi proyek; validasi dan dokumentasi. & 4 \\
12 & SCPMK0802-04802, SCPMK0607-04805 & Sprint 3: fitur lanjutan dan optimasi. & Modalitas: Blended Learning. Bentuk: Luring/Daring. Metode: Case Method, PBL, praktik proyek, dan presentasi. & 1 x 6 x 50' tatap muka; 1 x 6 x 50' tugas/praktik mandiri. & Mahasiswa mengerjakan aktivitas terarah untuk sprint 3 fitur lanjutan dan optimasi dan menunjukkan evidence hasil belajar. & Bentuk penilaian: Sprint 3 Fitur Lanjutan dan Optimasi. Mengacu rubrik RTM. & Ketepatan konsep; kualitas implementasi proyek; validasi dan dokumentasi. & 5 \\
13 & SCPMK0607-04805 & Sprint 3: deployment preparation. & Modalitas: Blended Learning. Bentuk: Luring/Daring. Metode: Case Method, PBL, praktik proyek, dan presentasi. & 1 x 6 x 50' tatap muka; 1 x 6 x 50' tugas/praktik mandiri. & Mahasiswa mengerjakan aktivitas terarah untuk sprint 3 deployment preparation dan menunjukkan evidence hasil belajar. & Bentuk penilaian: Sprint 3 Deployment Preparation. Mengacu rubrik RTM. & Ketepatan konsep; kualitas implementasi proyek; validasi dan dokumentasi. & 4 \\
14 & SCPMK0607-04805, SCPMK0804-04804 & Sprint 4: final integration dan UAT. & Modalitas: Blended Learning. Bentuk: Luring/Daring. Metode: Case Method, PBL, praktik proyek, dan presentasi. & 1 x 6 x 50' tatap muka; 1 x 6 x 50' tugas/praktik mandiri. & Mahasiswa mengerjakan aktivitas terarah untuk sprint 4 final integration dan UAT dan menunjukkan evidence hasil belajar. & Bentuk penilaian: Sprint 4 Final Integration dan UAT. Mengacu rubrik RTM. & Ketepatan konsep; kualitas implementasi proyek; validasi dan dokumentasi. & 5 \\
15 & SCPMK0804-04804, SCPMK0302-04807 & Penyusunan laporan akhir dan konsolidasi hasil kerja tim. & Modalitas: Blended Learning. Bentuk: Luring/Daring. Metode: Case Method, PBL, praktik proyek, dan presentasi. & 1 x 6 x 50' tatap muka; 1 x 6 x 50' tugas/praktik mandiri. & Mahasiswa mengerjakan aktivitas terarah untuk penyusunan laporan akhir, termasuk mengelola dinamika tim dan menyusun penyampaian hasil kerja tim, dan menunjukkan evidence hasil belajar. & Bentuk penilaian: Penyusunan Laporan Akhir. Mengacu rubrik RTM. & Ketepatan konsep; kualitas implementasi proyek; validasi dan dokumentasi; koordinasi tim dalam penyusunan laporan. & 5 \\
16 & SCPMK0804-04804 & Persiapan defense dan latihan presentasi. & Modalitas: Blended Learning. Bentuk: Luring/Daring. Metode: Case Method, PBL, praktik proyek, dan presentasi. & 1 x 6 x 50' tatap muka; 1 x 6 x 50' tugas/praktik mandiri. & Mahasiswa mengerjakan aktivitas terarah untuk persiapan defense dan latihan presentasi dan menunjukkan evidence hasil belajar. & Bentuk penilaian: Persiapan Defense dan Latihan Presentasi. Mengacu rubrik RTM. & Ketepatan konsep; kualitas implementasi proyek; validasi dan dokumentasi. & 4 \\
17 & SCPMK0803-04803, SCPMK0804-04804, SCPMK0302-04807 & UAS Defense proyek akhir dan demo, termasuk penyampaian hasil kerja tim kepada pemangku kepentingan. & Modalitas: Blended Learning. Bentuk: Luring/Daring. Metode: Case Method, PBL, praktik proyek, dan presentasi. & 1 x 6 x 50' tatap muka; 1 x 6 x 50' tugas/praktik mandiri. & Mahasiswa mengerjakan aktivitas terarah untuk UAS defense proyek akhir dan demo, termasuk penyampaian hasil kerja tim kepada pemangku kepentingan, dan menunjukkan evidence hasil belajar. & Bentuk penilaian: UAS Defense Proyek Akhir dan Demo. Mengacu rubrik RTM. & Ketepatan konsep; kualitas implementasi proyek; validasi dan dokumentasi; profesionalisme penyampaian hasil kepada pemangku kepentingan. & 20 \\
\end{longtblr}
